\section{Martingales}

\subsection{Conditional expectation}

\begin{de}[Conditioning on an event]
Let $(\Omega,\mathcal F,P)$ be a probability space and let $B\in\mathcal F$
with $P(B)>0$. The conditional probability measure is
\[
P_B(A)=P(A\mid B)=\frac{P(A\cap B)}{P(B)},
\qquad A\in\mathcal F.
\]
A measurable random variable $X$ belongs to $L^1(P_B)$ precisely when
\[
\mathbb E[|X|\mathbf1_B]<\infty,
\]
and then
\[
\mathbb E_B[X]=\frac{\mathbb E[X\mathbf1_B]}{P(B)}.
\]
In particular, $L^1(P)\subseteq L^1(P_B)$, but equality need not hold.
\end{de}

\begin{de}[Conditioning on a countable partition]
Let $(A_n)_{n\geq1}$ be a measurable partition of $\Omega$ with
$P(A_n)>0$, and put $\mathcal G=\sigma(A_n:n\geq1)$. For $X\in L^1$,
\[
\mathbb E[X\mid\mathcal G]
 =\sum_{n=1}^{\infty}
   \frac{\mathbb E[X\mathbf1_{A_n}]}{P(A_n)}\mathbf1_{A_n}.
\]
\end{de}

\begin{thm}[Existence and uniqueness of conditional expectation]
Let $X\in L^1(\Omega,\mathcal F,P)$ and let
$\mathcal G\subseteq\mathcal F$ be a sub-$\sigma$-field. There is a
$\mathcal G$-measurable random variable, unique almost surely, denoted
$\mathbb E[X\mid\mathcal G]$, such that
\[
\mathbb E[\mathbb E[X\mid\mathcal G]\mathbf1_A]
 =\mathbb E[X\mathbf1_A]
\qquad(A\in\mathcal G).
\]
\end{thm}

\begin{proof}
The signed measure
$\nu(A)=\mathbb E[X\mathbf1_A]$ on $(\Omega,\mathcal G)$ is absolutely
continuous with respect to $P|_{\mathcal G}$. Its Radon--Nikodym derivative
is a conditional expectation. If $Y$ and $Z$ both satisfy the defining
identity, apply it to $\{Y>Z\}$ and $\{Z>Y\}$ to obtain $Y=Z$ almost surely.
\end{proof}

\begin{prop}[$L^2$ projection interpretation]
The subspace $L^2(\Omega,\mathcal G,P)$ is closed in
$L^2(\Omega,\mathcal F,P)$. For $X\in L^2$,
$\mathbb E[X\mid\mathcal G]$ is the orthogonal projection of $X$ onto this
subspace. Equivalently, it is the unique $Z\in L^2(\mathcal G)$ satisfying
\[
\mathbb E[(X-Z)Y]=0
\qquad(Y\in L^2(\mathcal G)).
\]
\end{prop}

\begin{proof}
If $Y_n\in L^2(\mathcal G)$ and $Y_n\to Y$ in $L^2$, a subsequence converges
almost surely. Its pointwise limit is $\mathcal G$-measurable and is a version
of $Y$, so the subspace is closed. The characterization of orthogonal
projection gives the second assertion, and taking bounded
$\mathcal G$-measurable test functions identifies the projection with the
conditional expectation defined above.
\end{proof}

\begin{thm}[Basic properties]
Let $X,Y\in L^1$ and let $\mathcal G\subseteq\mathcal F$.
\begin{enumerate}[label=(\roman*)]
  \item Conditional expectation is linear.
  \item If $X\leq Y$ almost surely, then
  $\mathbb E[X\mid\mathcal G]\leq\mathbb E[Y\mid\mathcal G]$ almost surely.
  \item
  $|\mathbb E[X\mid\mathcal G]|
   \leq\mathbb E[|X|\mid\mathcal G]$ almost surely.
  \item If $Z$ is bounded and $\mathcal G$-measurable, then
  \[
  \mathbb E[ZX\mid\mathcal G]
   =Z\mathbb E[X\mid\mathcal G].
  \]
  The same holds whenever both sides are integrable.
  \item If $\mathcal G_1\subseteq\mathcal G_2\subseteq\mathcal F$, then
  \[
  \mathbb E[\mathbb E[X\mid\mathcal G_2]\mid\mathcal G_1]
   =\mathbb E[X\mid\mathcal G_1]
  \]
  and
  \[
  \mathbb E[\mathbb E[X\mid\mathcal G_1]\mid\mathcal G_2]
   =\mathbb E[X\mid\mathcal G_1].
  \]
  \item If $X$ is $\mathcal G$-measurable, then
  $\mathbb E[X\mid\mathcal G]=X$ almost surely.
  \item $\mathbb E[\mathbb E[X\mid\mathcal G]]=\mathbb E X$.
\end{enumerate}
\end{thm}

\begin{proof}
Linearity and the tower identities follow from the defining integral
identity and uniqueness. For monotonicity, let
$A=\{\mathbb E[X\mid\mathcal G]>
      \mathbb E[Y\mid\mathcal G]\}\in\mathcal G$; integration over $A$
forces $P(A)=0$. Applying monotonicity to $-\lvert X\rvert\leq X\leq
\lvert X\rvert$ gives (iii). For (iv), test both sides against indicators of
sets in $\mathcal G$. The remaining assertions are immediate consequences.
\end{proof}

\begin{de}[Nonnegative conditional expectation]
If $X\geq0$ is measurable, define
\[
\mathbb E[X\mid\mathcal G]
 =\lim_{n\to\infty}\mathbb E[X\wedge n\mid\mathcal G].
\]
The limit is $\mathcal G$-measurable, may equal $+\infty$, and satisfies the
defining integral identity for every $A\in\mathcal G$.
\end{de}

\begin{thm}[Conditional convergence theorems]
Let $\mathcal G\subseteq\mathcal F$.
\begin{enumerate}[label=(\roman*)]
  \item If $0\leq X_n\uparrow X$, then
  \[
  \mathbb E[X_n\mid\mathcal G]
   \uparrow\mathbb E[X\mid\mathcal G]
  \quad\text{almost surely}.
  \]
  \item If $X_n\geq0$, then
  \[
  \mathbb E[\liminf_nX_n\mid\mathcal G]
   \leq\liminf_n\mathbb E[X_n\mid\mathcal G]
  \quad\text{almost surely}.
  \]
  \item If $X_n\to X$ almost surely and $|X_n|\leq Y\in L^1$, then
  \[
  \mathbb E[X_n\mid\mathcal G]\to
  \mathbb E[X\mid\mathcal G]
  \quad\text{almost surely and in }L^1.
  \]
\end{enumerate}
\end{thm}

\begin{proof}
For (i), the increasing limit has the correct integral over every
$A\in\mathcal G$ by the ordinary monotone convergence theorem. Part (ii)
follows by applying (i) to
$Y_k=\inf_{n\geq k}X_n$. For (iii), apply conditional Fatou to
$Y\pm X_n$ to identify the almost sure limit. Moreover,
\[
\|\mathbb E[X_n-X\mid\mathcal G]\|_1
 \leq\mathbb E|X_n-X|\to0
\]
by ordinary dominated convergence.
\end{proof}

\begin{lem}[Countable supporting-line representation]
If $\varphi:\mathbb R\to\mathbb R$ is finite and convex, there is a countable
family of affine functions $(\ell_j)$ such that
$\ell_j\leq\varphi$ on $\mathbb R$ and
\[
\varphi(x)=\sup_{j\geq1}\ell_j(x).
\]
\end{lem}

\begin{proof}
If $\varphi$ is affine, one line suffices. Otherwise, for each rational $q$
choose a supporting line at $q$. Convexity makes subgradients locally bounded,
so supporting lines based at rationals approaching $x$ take values at $x$
that approach $\varphi(x)$. The countable collection therefore has the
stated supremum.
\end{proof}

\begin{thm}[Conditional Jensen inequality]
Let $X\in L^1$, let $\varphi:\mathbb R\to\mathbb R$ be convex, and suppose
$\varphi(X)\in L^1$. Then
\[
\varphi(\mathbb E[X\mid\mathcal G])
 \leq\mathbb E[\varphi(X)\mid\mathcal G]
\quad\text{almost surely}.
\]
In particular, for $1\leq p<\infty$,
\[
\|\mathbb E[X\mid\mathcal G]\|_p\leq\|X\|_p.
\]
\end{thm}

\begin{proof}
For every affine minorant $\ell_j(x)=a_jx+b_j$,
\[
\ell_j(\mathbb E[X\mid\mathcal G])
 =\mathbb E[\ell_j(X)\mid\mathcal G]
 \leq\mathbb E[\varphi(X)\mid\mathcal G].
\]
Take the countable supremum. An affine lower bound also shows that the
negative part of the left-hand side is integrable. The contraction follows
from $\varphi(x)=|x|^p$ and integration.
\end{proof}

\begin{prop}[Conditioning and independence]
Let $X:(\Omega,\mathcal F)\to(S,\mathcal S)$ and let
$\mathcal G\subseteq\mathcal F$. If $\sigma(X)$ and $\mathcal G$ are
independent, then for every integrable $f(X)$,
\[
\mathbb E[f(X)\mid\mathcal G]=\mathbb E f(X)
\quad\text{almost surely}.
\]
Conversely, if this holds for every bounded measurable $f$, then
$\sigma(X)$ and $\mathcal G$ are independent.
\end{prop}

\begin{proof}
For $A\in\mathcal G$, independence gives
$\mathbb E[f(X)\mathbf1_A]=\mathbb Ef(X)P(A)$, which is the defining identity.
Conversely, take $f=\mathbf1_B$ and test against $A\in\mathcal G$.
\end{proof}

\begin{de}[Conditioning on a random variable]
For $Y\in L^1$ and a random variable $X$, write
\[
\mathbb E[Y\mid X]
 :=\mathbb E[Y\mid\sigma(X)].
\]
If the state space of $X$ is standard Borel, the Doob--Dynkin lemma gives a
measurable function $g$ such that
$\mathbb E[Y\mid X]=g(X)$ almost surely.
\end{de}

\begin{prop}[Independent-parameter substitution]
Let $X$ be $\mathcal G$-measurable, let $Y$ be independent of $\mathcal G$,
and let $h:S\times T\to\mathbb R$ be measurable with
$h(X,Y)\in L^1$. Write $\nu=\mathcal L(Y)$ and define
\[
g(x)=
\begin{cases}
 \displaystyle\int_T h(x,y)\,\nu(\dd y),
   &\displaystyle\int_T|h(x,y)|\,\nu(\dd y)<\infty,\\[6pt]
 0,&\text{otherwise}.
\end{cases}
\]
Then $g$ is measurable and
\[
\mathbb E[h(X,Y)\mid\mathcal G]=g(X)
\quad\text{almost surely}.
\]
\end{prop}

\begin{proof}
For bounded nonnegative $h$, the claim is immediate for product indicators
$h(x,y)=u(x)v(y)$ and extends by the functional monotone-class theorem. The
general nonnegative case follows by monotone convergence. Applying this to
$|h|$ shows that the section integral is finite at $X$ almost surely; applying
it to $h^+$ and $h^-$ then proves the integrable case and the measurability of
$g$.
\end{proof}

\begin{eg}[Conditional density]
Suppose $(X,Z)$ has joint density $f_{X,Z}$ and let
$f_Z(z)=\int f_{X,Z}(x,z)\,\dd x$. Define
\[
f_{X\mid Z}(x\mid z)=
\begin{cases}
 f_{X,Z}(x,z)/f_Z(z),&f_Z(z)>0,\\
 0,&f_Z(z)=0.
\end{cases}
\]
If $h$ is Borel and $\mathbb E|h(X)|<\infty$, then
\[
\mathbb E[h(X)\mid Z]
 =g(Z),
\qquad
g(z)=\int h(x)f_{X\mid Z}(x\mid z)\,\dd x.
\]
\end{eg}

\begin{proof}
The function $g(Z)$ is $\sigma(Z)$-measurable. For a Borel set $B$,
Fubini's theorem gives
\begin{align*}
\mathbb E[g(Z)\mathbf1_{\{Z\in B\}}]
&=\int_B\int h(x)f_{X,Z}(x,z)\,\dd x\dd z\\
&=\mathbb E[h(X)\mathbf1_{\{Z\in B\}}].
\end{align*}
\end{proof}

\subsection{Basics of martingales}

\begin{de}[Stochastic process and its law]
A stochastic process with index set $I$ and state space $(E,\mathcal E)$ is a
family $X=(X_t)_{t\in I}$ of $E$-valued random variables. Its law is the
pushforward of $P$ under the path map
\[
\omega\longmapsto(X_t(\omega))_{t\in I}
\]
on $(E^I,\mathcal E^{\otimes I})$. For Brownian motion with continuous paths,
this law is usually viewed on $C([0,\infty))$ and is called Wiener measure.
\end{de}

\begin{de}[Common process properties]
A real-valued process has independent increments if increments over disjoint
time intervals are independent. It is Gaussian if every finite-dimensional
marginal is multivariate normal. It is stationary if
$(X_{t+s})_t$ and $(X_t)_t$ have the same finite-dimensional distributions
for every admissible shift $s$. It has stationary increments if
\[
(X_{t_1+s}-X_s,\ldots,X_{t_m+s}-X_s)
\stackrel d=
(X_{t_1}-X_0,\ldots,X_{t_m}-X_0)
\]
for all admissible $s,t_1,\ldots,t_m$.
\end{de}

From this point onward, time is discrete: $n\in\mathbb N_0$.

\begin{de}[Filtration and adapted process]
A filtration is an increasing family
$\mathbb F=(\mathcal F_n)_{n\geq0}$ of sub-$\sigma$-fields of $\mathcal F$.
A process $X=(X_n)$ is adapted if $X_n$ is $\mathcal F_n$-measurable. Its
natural filtration is
$\mathcal F_n^X=\sigma(X_0,\ldots,X_n)$.
\end{de}

\begin{de}[Martingales]
An adapted integrable process $X=(X_n)$ is a martingale, submartingale, or
supermartingale according as, for $m\leq n$,
\[
\mathbb E[X_n\mid\mathcal F_m]
 =X_m,
\qquad
\mathbb E[X_n\mid\mathcal F_m]
 \geq X_m,
\qquad\text{or}\qquad
\mathbb E[X_n\mid\mathcal F_m]
 \leq X_m.
\]
\end{de}

\begin{de}[Stopping time]
A map $\tau:\Omega\to\mathbb N_0\cup\{\infty\}$ is a stopping time if
\[
\{\tau\leq n\}\in\mathcal F_n
\qquad(n\geq0).
\]
Equivalently, $\{\tau=n\}\in\mathcal F_n$ for every $n$.
\end{de}

\begin{prop}[Operations on stopping times]
If $\sigma$ and $\tau$ are stopping times, then
$\sigma\wedge\tau$, $\sigma\vee\tau$, and $\sigma+\tau$ are stopping times.
For every deterministic $r\in\mathbb N_0$, $\tau+r$ is a stopping time.
In general, $\tau-r$ need not be a stopping time.
\end{prop}

\begin{proof}
Use
\[
\{\sigma\wedge\tau\leq n\}
 =\{\sigma\leq n\}\cup\{\tau\leq n\},
\qquad
\{\sigma\vee\tau\leq n\}
 =\{\sigma\leq n\}\cap\{\tau\leq n\},
\]
and
\[
\{\sigma+\tau\leq n\}
 =\bigcup_{k=0}^{n}\{\sigma=k\}\cap\{\tau\leq n-k\}.
\]
Every term in the final finite union belongs to $\mathcal F_n$.
\end{proof}

\begin{de}[Stopping-time $\sigma$-field]
For a stopping time $\tau$, define
\[
\mathcal F_\tau
 =\{A\in\mathcal F:A\cap\{\tau\leq n\}\in\mathcal F_n
   \text{ for every }n\}.
\]
\end{de}

\begin{prop}
If $\sigma\leq\tau$, then
$\mathcal F_\sigma\subseteq\mathcal F_\tau$. Moreover,
\[
\mathcal F_{\sigma\wedge\tau}
 =\mathcal F_\sigma\cap\mathcal F_\tau.
\]
If $X$ is adapted and $\tau<\infty$ almost surely, then
$X_\tau$ is $\mathcal F_\tau$-measurable.
\end{prop}

\begin{proof}
For the first inclusion, if $A\in\mathcal F_\sigma$, then
\[
A\cap\{\tau\leq n\}
 =(A\cap\{\sigma\leq n\})\cap\{\tau\leq n\}\in\mathcal F_n.
\]
The equality for the minimum follows from
$\{\sigma\wedge\tau\leq n\}=
\{\sigma\leq n\}\cup\{\tau\leq n\}$. Finally, for Borel $B$,
\[
\{X_\tau\in B\}\cap\{\tau\leq n\}
 =\bigcup_{k=0}^{n}\{\tau=k\}\cap\{X_k\in B\}\in\mathcal F_n.
\]
\end{proof}

\begin{prop}[Elementary closure properties]
\begin{enumerate}[label=(\roman*)]
  \item $X$ is a supermartingale if and only if $-X$ is a submartingale.
  \item Linear combinations of martingales are martingales.
  \item Nonnegative linear combinations of submartingales, respectively
  supermartingales, stay in the same class.
  \item The maximum of two submartingales is a submartingale, and the minimum
  of two supermartingales is a supermartingale.
  \item If $X$ is a supermartingale and
  $\mathbb EX_N\geq\mathbb EX_0$ for some $N$, then $(X_n)_{0\leq n\leq N}$
  is a martingale.
\end{enumerate}
\end{prop}

\begin{proof}
Only (iv) needs a calculation. Using
$x\vee y=(x+y+|x-y|)/2$ and conditional Jensen for the absolute-value
function gives
\[
\mathbb E[X_n\vee Y_n\mid\mathcal F_m]
 \geq\mathbb E[X_n\mid\mathcal F_m]
       \vee\mathbb E[Y_n\mid\mathcal F_m]
 \geq X_m\vee Y_m.
\]
The supermartingale statement follows by negation. For (v), expectations of a
supermartingale are nonincreasing. Equality at the endpoints forces equality
at every intermediate step, and the nonpositive conditional drifts must
vanish.
\end{proof}

\begin{thm}[Convex transforms]
Let $X$ be a martingale and let $\varphi:\mathbb R\to\mathbb R$ be convex.
If $\varphi(X_n)$ is integrable for every $n$, then
$(\varphi(X_n))$ is a submartingale. More generally, if $X$ is a
submartingale and $\varphi$ is increasing and convex, then
$(\varphi(X_n))$ is a submartingale whenever it is integrable.
Consequently,
\[
|X_n|^p\quad(p\geq1),
\qquad
(X_n-a)^+,
\]
are submartingales in the applicable cases, and if $X$ is a supermartingale,
then $X_n\wedge a$ is a supermartingale.
\end{thm}

\begin{proof}
Conditional Jensen gives
\[
\mathbb E[\varphi(X_n)\mid\mathcal F_m]
 \geq\varphi(\mathbb E[X_n\mid\mathcal F_m]).
\]
For a martingale this equals $\varphi(X_m)$; for a submartingale, monotonicity
of $\varphi$ makes it at least $\varphi(X_m)$. Finally,
$X\wedge a=a-(a-X)^+$ and $-X$ is a submartingale when $X$ is a
supermartingale.
\end{proof}

\subsection{Optional sampling and stochastic integration}

\begin{de}[Predictable process]
A discrete-time process $H=(H_n)_{n\geq1}$ is predictable if
$H_n$ is $\mathcal F_{n-1}$-measurable for every $n$.
\end{de}

\begin{thm}[Doob decomposition]
Every adapted integrable process $X=(X_n)$ has a unique decomposition
\[
X_n=M_n+A_n,
\]
where $M$ is a martingale, $A_0=0$, and $A$ is predictable. Explicitly,
\[
A_n=\sum_{k=1}^{n}
 \bigl(\mathbb E[X_k\mid\mathcal F_{k-1}]-X_{k-1}\bigr),
\qquad
M_n=X_n-A_n.
\]
The process $X$ is a submartingale if and only if $A$ is increasing, and it
is a supermartingale if and only if $A$ is decreasing.
\end{thm}

\begin{proof}
The increment
$M_k-M_{k-1}=X_k-\mathbb E[X_k\mid\mathcal F_{k-1}]$ has conditional mean
zero, so $M$ is a martingale. The assertions about monotonicity follow
straight from the increments of $A$. If
$X=M+A=M'+A'$, then $M-M'=A'-A$ is a predictable martingale starting at zero.
Its increments are both $\mathcal F_{n-1}$-measurable and have conditional
mean zero, hence vanish almost surely. This proves uniqueness.
\end{proof}

\begin{de}[Predictable quadratic variation]
For a square-integrable martingale $M$, its predictable quadratic variation
is
\[
\langle M\rangle_n
 =\sum_{k=1}^{n}
   \mathbb E[(M_k-M_{k-1})^2\mid\mathcal F_{k-1}].
\]
It is the unique predictable increasing process starting at zero such that
$M_n^2-\langle M\rangle_n$ is a martingale.
\end{de}

\begin{prop}
For every square-integrable martingale $M$,
\[
\mathbb E\langle M\rangle_n
 =\mathbb E[(M_n-M_0)^2]
\]
and
\[
\mathbb E M_n^2
 =\mathbb E M_0^2+\mathbb E\langle M\rangle_n.
\]
\end{prop}

\begin{proof}
Martingale differences are pairwise orthogonal in $L^2$ and are orthogonal
to $M_0$. Summing their second moments gives both identities.
\end{proof}

\begin{thm}[Optional sampling for bounded stopping times]
Let $X$ be a supermartingale and let $\sigma\leq\tau\leq N$ be stopping times.
Then
\[
\mathbb E[X_\tau\mid\mathcal F_\sigma]\leq X_\sigma
\quad\text{almost surely},
\]
and therefore $\mathbb EX_\tau\leq\mathbb EX_\sigma$. For a martingale,
equality holds. Conversely, an adapted integrable process $X$ is a martingale
if and only if
\[
\mathbb EX_\tau=\mathbb EX_0
\]
for every bounded stopping time $\tau$.
\end{thm}

\begin{proof}
For $A\in\mathcal F_\sigma$,
\[
X_\tau-X_\sigma
 =\sum_{k=1}^{N}
  \mathbf1_{\{\sigma<k\leq\tau\}}(X_k-X_{k-1}).
\]
The factor
$\mathbf1_A\mathbf1_{\{\sigma<k\leq\tau\}}$ is
$\mathcal F_{k-1}$-measurable. Taking expectations shows
$\mathbb E[(X_\tau-X_\sigma)\mathbf1_A]\leq0$, which is the conditional
inequality. For the converse, fix $m<n$ and $A\in\mathcal F_m$. Apply the
assumption to the stopping time equal to $m$ on $A$ and $n$ on $A^c$, and
compare it with the deterministic stopping time $n$. This gives
$\mathbb E[X_m\mathbf1_A]=\mathbb E[X_n\mathbf1_A]$.
\end{proof}

\begin{cor}[Nonnegative supermartingales]
Let $X$ be a nonnegative supermartingale and let
$\sigma\leq\tau<\infty$ almost surely. Then $X_\sigma,X_\tau\in L^1$ and
\[
\mathbb E[X_\tau\mid\mathcal F_\sigma]\leq X_\sigma.
\]
\end{cor}

\begin{proof}
Optional sampling applied to $\tau\wedge n$ gives
$\mathbb E X_{\tau\wedge n}\leq\mathbb E X_0$. Fatou's lemma gives
$\mathbb E X_\tau\leq\mathbb E X_0$, and similarly for $\sigma$. For
$A\in\mathcal F_\sigma$, apply bounded optional sampling to
$\sigma\wedge n$ and $\tau\wedge n$ and integrate over
$A\cap\{\sigma\leq n\}\in\mathcal F_{\sigma\wedge n}$. This gives
\[
\mathbb E[X_{\tau\wedge n}\mathbf1_{A\cap\{\sigma\leq n\}}]
 \leq
\mathbb E[X_\sigma\mathbf1_{A\cap\{\sigma\leq n\}}].
\]
Fatou's lemma on the left and monotone convergence on the right yield
$\mathbb E[X_\tau\mathbf1_A]\leq\mathbb E[X_\sigma\mathbf1_A]$.
\end{proof}

\begin{de}[Stopped process]
For a stopping time $\tau$, define
\[
X_n^\tau=X_{n\wedge\tau},
\qquad
\mathcal F_n^\tau=\mathcal F_{n\wedge\tau}.
\]
\end{de}

\begin{thm}[Stopping preserves martingales]
If $X$ is a martingale, submartingale, or supermartingale, then $X^\tau$ has
the same property with respect to both $\mathbb F$ and the stopped filtration
$\mathbb F^\tau$.
\end{thm}

\begin{proof}
For a submartingale,
\[
X_n^\tau-X_{n-1}^\tau
 =\mathbf1_{\{\tau\geq n\}}(X_n-X_{n-1}),
\]
and the indicator is $\mathcal F_{n-1}$-measurable. Conditional expectation
therefore gives a nonnegative drift. Conditioning once more on the smaller
$\sigma$-field $\mathcal F_{n-1}^\tau$ proves the assertion for the stopped
filtration. The other cases follow similarly or by negation.
\end{proof}

\begin{de}[Discrete stochastic integral]
Let $X$ be adapted and let $H$ be predictable. Whenever the summands are
integrable, define
\[
(H\mathbin{\cdot}X)_n
 =\sum_{k=1}^{n}H_k(X_k-X_{k-1}).
\]
\end{de}

\begin{thm}[Stability under bounded predictable transforms]
Let $X$ be adapted and integrable.
\begin{enumerate}[label=(\roman*)]
  \item $X$ is a martingale if and only if $H\mathbin{\cdot}X$ is a
  martingale for every bounded predictable process $H$.
  \item $X$ is a submartingale, respectively supermartingale, if and only if
  $H\mathbin{\cdot}X$ has the same property for every bounded nonnegative
  predictable $H$.
\end{enumerate}
\end{thm}

\begin{proof}
If $X$ is a martingale, then
\[
\mathbb E[H_n(X_n-X_{n-1})\mid\mathcal F_{n-1}]=0.
\]
The sign version is identical. Conversely, choose
$H_k=\mathbf1_{\{k=n\}}$ to isolate each one-step increment.
\end{proof}

\begin{remark}
Taking $H_k=\mathbf1_{\{\tau\geq k\}}$ gives
\[
X_{n\wedge\tau}=X_0+(H\mathbin{\cdot}X)_n,
\]
which is another proof that stopping preserves the martingale class.
\end{remark}

\subsection{Convergence theorems}

\subsubsection{Almost sure convergence}

\begin{de}[Upcrossings]
For a real sequence $x=(x_n)$ and $a<b$, define successively
\[
S_1=\inf\{n:x_n\leq a\},
\qquad
T_1=\inf\{n\geq S_1:x_n\geq b\},
\]
and
\[
S_{k+1}=\inf\{n\geq T_k:x_n\leq a\},
\qquad
T_{k+1}=\inf\{n\geq S_{k+1}:x_n\geq b\},
\]
with $\inf\varnothing=\infty$. The number of completed upcrossings by time
$n$ is
\[
U_n[a,b](x)=\sum_{k\geq1}\mathbf1_{\{T_k\leq n\}},
\]
and $U_\infty[a,b]=\lim_nU_n[a,b]$.
\end{de}

\begin{lem}[Deterministic upcrossing criterion]
A real sequence $(x_n)$ converges in $[-\infty,\infty]$ if and only if
$U_\infty[a,b](x)<\infty$ for every pair of rational numbers $a<b$.
\end{lem}

\begin{proof}
A convergent sequence can cross a fixed interval only finitely many times.
Conversely, if $\liminf x_n<\limsup x_n$, choose rationals
$\liminf x_n<a<b<\limsup x_n$. The sequence then completes infinitely many
upcrossings of $[a,b]$.
\end{proof}

\begin{thm}[Doob upcrossing inequality]
If $X=(X_n)$ is a submartingale, then for every $a<b$ and $n$,
\[
(b-a)\mathbb E U_n[a,b](X)
 \leq
\mathbb E[(X_n-a)^+]-\mathbb E[(X_0-a)^+].
\]
\end{thm}

\begin{proof}
Let $Y_n=(X_n-a)^+$, a submartingale. Define the predictable process
\[
H_j=\sum_{k\geq1}\mathbf1_{\{S_k<j\leq T_k\}},
\]
which takes values in $\{0,1\}$. The gain from every completed transaction is
at least $b-a$, and an unfinished transaction has nonnegative value, so
\[
(H\mathbin{\cdot}Y)_n\geq(b-a)U_n[a,b].
\]
The process $K=1-H$ is also nonnegative and predictable, hence
$K\mathbin{\cdot}Y$ is a submartingale starting at zero. Since
\[
(H\mathbin{\cdot}Y)_n+(K\mathbin{\cdot}Y)_n=Y_n-Y_0,
\]
we have
$\mathbb E(H\mathbin{\cdot}Y)_n\leq\mathbb E(Y_n-Y_0)$, proving the result.
\end{proof}

\begin{thm}[Martingale convergence theorem]
If a submartingale or supermartingale $(X_n)$ is bounded in $L^1$, then there
is $X_\infty\in L^1$ such that
\[
X_n\longrightarrow X_\infty
\qquad\text{almost surely}.
\]
\end{thm}

\begin{proof}
It is enough to treat a submartingale. For rational $a<b$, the upcrossing
inequality and the $L^1$ bound give
$\sup_n\mathbb EU_n[a,b]<\infty$. Monotone convergence implies
$\mathbb EU_\infty[a,b]<\infty$, so the number of upcrossings is finite almost
surely. Intersecting over the countable set of rational pairs and applying the
deterministic criterion gives an almost sure limit in the extended reals.
Fatou's lemma gives
\[
\mathbb E|X_\infty|
 \leq\liminf_n\mathbb E|X_n|<\infty,
\]
so the limit is finite almost surely.
\end{proof}

\begin{cor}
Every nonnegative supermartingale converges almost surely to some
$X_\infty\in L^1$, and
\[
X_n\geq\mathbb E[X_\infty\mid\mathcal F_n]
\qquad(n\geq0).
\]
\end{cor}

\begin{proof}
The family is $L^1$-bounded because
$\mathbb EX_n\leq\mathbb EX_0$. For $m\geq n$,
$X_n\geq\mathbb E[X_m\mid\mathcal F_n]$. Conditional Fatou gives
\[
\mathbb E[X_\infty\mid\mathcal F_n]
 \leq\liminf_{m\to\infty}
      \mathbb E[X_m\mid\mathcal F_n]
 \leq X_n.
\]
\end{proof}

\subsubsection{Uniform integrability and \texorpdfstring{$L^1$}{L1} convergence}

\begin{de}[Uniform integrability]
A family $\mathcal X\subset L^1(P)$ is uniformly integrable if
\[
\lim_{K\to\infty}
 \sup_{X\in\mathcal X}
 \mathbb E[|X|\mathbf1_{\{|X|>K\}}]=0.
\]
\end{de}

\begin{thm}[de la Vall\'ee--Poussin criterion]
A family $\mathcal X\subset L^1(P)$ is uniformly integrable if and only if
there is an increasing convex function
$\Phi:[0,\infty)\to[0,\infty)$ such that
\[
\frac{\Phi(x)}x\longrightarrow\infty
\quad\text{and}\quad
\sup_{X\in\mathcal X}\mathbb E\Phi(|X|)<\infty.
\]
\end{thm}

\begin{proof}
If such $\Phi$ exists, let
$c_K=\inf_{x\geq K}\Phi(x)/x\to\infty$. Then
\[
\mathbb E[|X|\mathbf1_{\{|X|\geq K\}}]
 \leq c_K^{-1}\mathbb E\Phi(|X|),
\]
uniformly in $X$. Conversely, choose $a_n\uparrow\infty$ so that
\[
\sup_{X\in\mathcal X}
\mathbb E[(|X|-a_n)^+]<2^{-n},
\]
and put
$\Phi(x)=\sum_{n\geq1}(x-a_n)^+$. This is increasing and convex, its ratio to
$x$ tends to infinity, and monotone convergence gives
$\sup_X\mathbb E\Phi(|X|)\leq1$.
\end{proof}

\begin{cor}
Every $L^p$-bounded family with $p>1$ is uniformly integrable. In particular,
a family with uniformly bounded means in absolute value and uniformly bounded
variances is uniformly integrable.
\end{cor}

\begin{proof}
Use $\Phi(x)=x^p$. In the second assertion,
$\mathbb EX^2=(\mathbb EX)^2+\operatorname{Var}(X)$ is uniformly bounded.
\end{proof}

\begin{prop}[Elementary stability]
Finite subsets of $L^1$ are uniformly integrable. If $\mathcal X$ and
$\mathcal Y$ are uniformly integrable, then so are
$\mathcal X+\mathcal Y$, $\mathcal X-\mathcal Y$, and
$\{|X|:X\in\mathcal X\}$. If each $Y$ in a family $\mathcal Y$ satisfies
$|Y|\leq|X|$ for some $X\in\mathcal X$, then uniform integrability of
$\mathcal X$ implies that of $\mathcal Y$.
\end{prop}

\begin{thm}[Vitali convergence theorem]
For random variables $(X_n)$ and $X$, the following are equivalent:
\begin{enumerate}[label=(\roman*)]
  \item $X_n\to X$ in $L^1$;
  \item $(X_n)$ is Cauchy in $L^1$ and $X_n\to X$ in probability;
  \item $(X_n)$ is uniformly integrable and $X_n\to X$ in probability.
\end{enumerate}
\end{thm}

\begin{proof}
Clearly (i) implies (ii). If (ii) holds, completeness of $L^1$ gives an
$L^1$ limit $Y$. Then $X_n\to Y$ in probability, so uniqueness of limits in
probability gives $X=Y$ almost surely and proves (i).

If (i) holds, convergence in probability follows. To prove uniform
integrability, choose $N$ so that
$\|X_n-X\|_1<\varepsilon$ for $n\geq N$. For the tail of the sequence,
\[
\mathbb E[|X_n|\mathbf1_{\{|X_n|>K\}}]
 \leq\varepsilon+
 \mathbb E[|X|\mathbf1_{\{|X_n|>K\}}].
\]
The expectations $\mathbb E|X_n|$ are uniformly bounded for $n\geq N$, so
Markov's inequality makes the probabilities of the displayed events uniformly
small as $K\to\infty$. Absolute continuity of the integral of $|X|$ then
controls the second term; the finitely many earlier variables cause no
problem. Thus (i) implies (iii).

Conversely, assume (iii). Uniform integrability implies
$\sup_n\mathbb E|X_n|<\infty$. An almost surely convergent subsequence and
Fatou's lemma show that $X\in L^1$. For $K>0$,
\begin{align*}
\mathbb E|X_n-X|
&\leq \mathbb E(|X_n-X|\wedge K)\\
&\quad+2\mathbb E[|X_n|\mathbf1_{\{|X_n|>K/2\}}]
     +2\mathbb E[|X|\mathbf1_{\{|X|>K/2\}}].
\end{align*}
For fixed $K$, the first term tends to zero because the bounded variables
$|X_n-X|\wedge K$ converge to zero in probability. The last two terms are
uniformly small for large $K$, proving $L^1$ convergence.
\end{proof}

\begin{lem}[Conditional expectations of one variable are uniformly integrable]
For every $Z\in L^1$, the family
\[
\{\mathbb E[Z\mid\mathcal G]:\mathcal G\subseteq\mathcal F
  \text{ a sub-$\sigma$-field}\}
\]
is uniformly integrable.
\end{lem}

\begin{proof}
Put $Y=\mathbb E[Z\mid\mathcal G]$ and
$A=\{|Y|>K\}\in\mathcal G$. Then
\[
K P(A)\leq\mathbb E[|Y|\mathbf1_A]
 \leq\mathbb E[|Z|\mathbf1_A].
\]
Also $P(A)\leq\mathbb E|Z|/K$. Absolute continuity of the integral of $|Z|$
therefore makes the final term uniformly small as $K\to\infty$.
\end{proof}

\begin{thm}[$L^1$ martingale convergence]
For a martingale $(X_n)$, the following are equivalent:
\begin{enumerate}[label=(\roman*)]
  \item $X_n\to X_\infty$ almost surely and in $L^1$;
  \item there is $Z\in L^1$ such that
  $X_n=\mathbb E[Z\mid\mathcal F_n]$ for every $n$;
  \item the family $(X_n)$ is uniformly integrable.
\end{enumerate}
When these conditions hold, one may take $Z=X_\infty$, and
\[
X_n=\mathbb E[X_\infty\mid\mathcal F_n].
\]
\end{thm}

\begin{proof}
If (i) holds, then for $m\geq n$,
$X_n=\mathbb E[X_m\mid\mathcal F_n]$. Contractivity of conditional
expectation allows $m\to\infty$ in $L^1$, giving (ii) with $Z=X_\infty$.
Condition (ii) implies (iii) by the preceding lemma. If (iii) holds, the
family is $L^1$-bounded, so the martingale convergence theorem gives an almost
sure limit. Vitali's theorem upgrades the convergence to $L^1$.
\end{proof}

\begin{thm}[Optional sampling for uniformly integrable martingales]
Let $(X_n)$ be a uniformly integrable martingale with $L^1$ limit
$X_\infty$. For a stopping time $T$, define $X_T=X_\infty$ on
$\{T=\infty\}$. Then
\[
X_T=\mathbb E[X_\infty\mid\mathcal F_T].
\]
Consequently, $X_T\in L^1$, the family of all stopped values $(X_T)$ is
uniformly integrable, and if $S\leq T$, then
\[
X_S=\mathbb E[X_T\mid\mathcal F_S].
\]
\end{thm}

\begin{proof}
We have $X_n=\mathbb E[X_\infty\mid\mathcal F_n]$. Conditional Jensen gives
\[
\mathbb E|X_T|
 \leq\sum_{n\geq0}
 \mathbb E[|X_\infty|\mathbf1_{\{T=n\}}]
 +\mathbb E[|X_\infty|\mathbf1_{\{T=\infty\}}]
 =\mathbb E|X_\infty|.
\]
For $A\in\mathcal F_T$, decompose according to $\{T=n\}$ and
$\{T=\infty\}$ to obtain
$\mathbb E[X_T\mathbf1_A]=\mathbb E[X_\infty\mathbf1_A]$. This proves the
conditional-expectation identity. Uniform integrability follows from the
preceding lemma, and the final assertion follows from the tower property and
$\mathcal F_S\subseteq\mathcal F_T$.
\end{proof}

\begin{thm}[Optional sampling for nonnegative or uniformly integrable supermartingales]
Let $X$ be a supermartingale and suppose either
\begin{enumerate}[label=(\alph*)]
  \item $X_n\geq0$ for every $n$, or
  \item $(X_n)$ is uniformly integrable.
\end{enumerate}
If $S\leq T<\infty$ almost surely, then $X_S,X_T\in L^1$ and
\[
\mathbb E[X_T\mid\mathcal F_S]\leq X_S.
\]
\end{thm}

\begin{proof}
Case (a) was proved above. For case (b), write the Doob decomposition as
$X=M-A$, where $A$ is predictable, increasing, and $A_0=0$. Since
\[
\mathbb EA_n=\mathbb EX_0-\mathbb EX_n
\]
is bounded, $A_n\uparrow A_\infty$ in $L^1$. Thus $(A_n)$ is uniformly
integrable, and so is the martingale $M=X+A$. The martingale is therefore
closed, and all its stopped values are uniformly integrable; also
$0\leq A_T\leq A_\infty$. Hence $X_{T\wedge n}\to X_T$ and
$X_{S\wedge n}\to X_S$ in $L^1$. For $B\in\mathcal F_S$, integrate bounded
optional sampling over
$B\cap\{S\leq n\}\in\mathcal F_{S\wedge n}$ and then let $n\to\infty$.
The $L^1$ convergences and
$\mathbf1_{B\cap\{S\leq n\}}\to\mathbf1_B$ give
\[
\mathbb E[X_T\mathbf1_B]\leq\mathbb E[X_S\mathbf1_B],
\]
which is the asserted conditional inequality.
\end{proof}

\begin{eg}[Gambler's ruin]
Let $(S_n)$ be simple symmetric random walk with $S_0=0$, and let
$a<0<b$ be integers. Define
\[
\tau=\inf\{n\geq0:S_n\in\{a,b\}\}.
\]
Then $\tau<\infty$ almost surely and
\[
P(S_\tau=a)=\frac{b}{b-a},
\qquad
P(S_\tau=b)=\frac{-a}{b-a}.
\]
Indeed, from every interior state there is a probability at least
$2^{-(b-a)}$ of hitting a boundary within the next $b-a$ steps. Iterating over
blocks shows that $P(\tau>m(b-a))\to0$. The stopped martingale is bounded, so
optional sampling gives $0=\mathbb ES_\tau$, and the two displayed
probabilities follow from their sum being one.
\end{eg}

\begin{eg}[Expected exit time]
For the same stopping time,
$S_n^2-n$ is a martingale. Bounded optional sampling gives
\[
\mathbb E(\tau\wedge n)=\mathbb E S_{\tau\wedge n}^2
 \leq\max\{a^2,b^2\}.
\]
Monotone convergence therefore gives $\mathbb E\tau<\infty$, while dominated
convergence applies to the bounded stopped positions. Hence
\[
\mathbb E\tau
 =\mathbb E S_\tau^2
 =a^2\frac{b}{b-a}+b^2\frac{-a}{b-a}
 =|a|b.
\]
If $\tau_a$ is the first hitting time of a fixed level $a<0$, then the events
that $a$ is hit before $b$ increase to $\{\tau_a<\infty\}$ as
$b\to\infty$. Their probabilities tend to one by the gambler's-ruin formula,
so $\tau_a<\infty$ almost surely. Moreover, the exit times increase to
$\tau_a$, and their expectations $|a|b$ tend to infinity. Thus
$\mathbb E\tau_a=\infty$.
\end{eg}

\subsubsection{\texorpdfstring{$L^p$}{Lp} convergence for \texorpdfstring{$p>1$}{p greater than 1}}

\begin{lem}[Doob maximal inequality, weak form]
If $X$ is a submartingale, set
\[
X_n^*=\max_{0\leq k\leq n}X_k.
\]
Then, for every $\lambda>0$,
\[
\lambda P(X_n^*\geq\lambda)
 \leq\mathbb E[X_n\mathbf1_{\{X_n^*\geq\lambda\}}].
\]
\end{lem}

\begin{proof}
Let $\tau=\inf\{k:X_k\geq\lambda\}\wedge n$. Bounded optional sampling gives
$\mathbb EX_\tau\leq\mathbb EX_n$. On the event
$\{X_n^*<\lambda\}$, $X_\tau=X_n$, while on its complement
$X_\tau\geq\lambda$. Rearranging yields the inequality.
\end{proof}

\begin{thm}[Doob's $L^p$ inequalities]
Let $X$ be a martingale, or a nonnegative submartingale, and put
$|X|_n^*=\max_{k\leq n}|X_k|$.
\begin{enumerate}[label=(\roman*)]
  \item For $p\geq1$ and $\lambda>0$,
  \[
  \lambda^pP(|X|_n^*\geq\lambda)
   \leq\mathbb E|X_n|^p.
  \]
  \item For $p>1$,
  \[
  \mathbb E|X_n|^p
   \leq\mathbb E(|X|_n^*)^p
   \leq\left(\frac p{p-1}\right)^p\mathbb E|X_n|^p.
  \]
\end{enumerate}
\end{thm}

\begin{proof}
The process $|X|^p$ is a nonnegative submartingale, so the weak inequality
follows from the lemma. For the strong inequality, apply the weak form to
$|X|$ and use the layer-cake formula. For $K>0$,
\begin{align*}
\mathbb E[(|X|_n^*\wedge K)^p]
&=\int_0^Kp\lambda^{p-1}P(|X|_n^*\geq\lambda)\,\dd\lambda\\
&\leq p\mathbb E\!\left[
 |X_n|\int_0^{|X|_n^*\wedge K}\lambda^{p-2}\,\dd\lambda
 \right]\\
&=\frac p{p-1}
 \mathbb E[|X_n|(|X|_n^*\wedge K)^{p-1}].
\end{align*}
Hölder's inequality gives the stated constant. Let $K\to\infty$.
\end{proof}

\begin{thm}[$L^p$ martingale convergence]
Let $p>1$ and let $(X_n)$ be a martingale satisfying
\[
\sup_n\mathbb E|X_n|^p<\infty.
\]
With $\mathcal F_\infty=\sigma(\bigcup_n\mathcal F_n)$, there is
$X_\infty\in L^p(\mathcal F_\infty)$ such that
\[
X_n\to X_\infty
\quad\text{almost surely and in }L^p.
\]
Moreover, $(|X_n|^p)$ is uniformly integrable.
\end{thm}

\begin{proof}
The $L^p$ bound implies uniform integrability, so an almost sure and $L^1$
limit exists. Doob's inequality and monotone convergence give
\[
\mathbb E\sup_{n\geq0}|X_n|^p
 \leq\left(\frac p{p-1}\right)^p
      \sup_n\mathbb E|X_n|^p<\infty.
\]
Thus $|X_n-X_\infty|^p$ is dominated by an integrable multiple of
$\sup_k|X_k|^p$, and dominated convergence gives $L^p$ convergence. The same
dominator proves uniform integrability of $(|X_n|^p)$.
\end{proof}

\begin{cor}[Square-integrable criterion]
For a square-integrable martingale $M$, the following are equivalent:
\begin{enumerate}[label=(\roman*)]
  \item $\sup_n\mathbb EM_n^2<\infty$;
  \item $\sup_n\mathbb E\langle M\rangle_n<\infty$;
  \item $M_n$ converges in $L^2$;
  \item $M_n$ converges almost surely and in $L^2$.
\end{enumerate}
\end{cor}

\begin{proof}
The identity
$\mathbb EM_n^2=\mathbb EM_0^2+\mathbb E\langle M\rangle_n$ gives
(i)$\Leftrightarrow$(ii). The $L^2$ martingale convergence theorem gives
(i)$\Rightarrow$(iv), while (iv)$\Rightarrow$(iii)$\Rightarrow$(i) is
immediate.
\end{proof}

\begin{thm}[Almost sure convergence from finite quadratic variation]
If $M$ is square integrable and
\[
\sup_n\langle M\rangle_n<\infty
\quad\text{almost surely},
\]
then $M_n$ converges almost surely.
\end{thm}

\begin{proof}
Assume $M_0=0$. For $K>0$, set
\[
\tau_K=\inf\{n:\langle M\rangle_{n+1}>K\}.
\]
Predictability makes $\tau_K$ a stopping time, and
$\langle M^{\tau_K}\rangle_n\leq K$. The square-integrable criterion shows
that $M^{\tau_K}$ converges almost surely. Since
$P(\tau_K=\infty)\to1$ as $K\to\infty$, the original martingale converges on
a full-probability union of these events.
\end{proof}

\subsubsection{Applications of martingale convergence}

\begin{thm}[L\'evy's upward theorem]
Let $X\in L^1$ and let
$\mathcal F_n\uparrow\mathcal F_\infty:=
\sigma(\bigcup_n\mathcal F_n)$. Then
\[
\mathbb E[X\mid\mathcal F_n]
 \longrightarrow\mathbb E[X\mid\mathcal F_\infty]
\quad\text{almost surely and in }L^1.
\]
\end{thm}

\begin{proof}
The process $Y_n=\mathbb E[X\mid\mathcal F_n]$ is a closed, hence uniformly
integrable, martingale. Let $Y_\infty$ be its $L^1$ limit. For
$A\in\bigcup_n\mathcal F_n$,
$\mathbb E[Y_\infty\mathbf1_A]=\mathbb E[X\mathbf1_A]$. A
$\pi$--$\lambda$ argument extends the identity to all
$A\in\mathcal F_\infty$, identifying
$Y_\infty=\mathbb E[X\mid\mathcal F_\infty]$.
\end{proof}

\begin{cor}[L\'evy's zero--one form]
If $A\in\mathcal F_\infty$, then
\[
\mathbb E[\mathbf1_A\mid\mathcal F_n]\to\mathbf1_A
\quad\text{almost surely and in }L^1.
\]
\end{cor}

\begin{remark}
For independent $(X_n)$ and
$\mathcal F_n=\sigma(X_1,\ldots,X_n)$, a tail event $A$ is independent of
each $\mathcal F_n$, so
$\mathbb E[\mathbf1_A\mid\mathcal F_n]=P(A)$. L\'evy's theorem therefore
gives $\mathbf1_A=P(A)$ almost surely, recovering Kolmogorov's zero--one law.
\end{remark}

\subsection{Reverse martingales and exchangeability}

\begin{de}[Exchangeability]
Let $E$ be a standard Borel space. A sequence $(X_n)_{n\geq1}$ of $E$-valued
random variables is exchangeable if its joint law is invariant under every
finite permutation of the coordinates.
\end{de}

Let $S_n$ denote the permutations of $\mathbb N$ that fix every index larger
than $n$. On the canonical space $E^{\mathbb N}$, let $\mathcal E_n'$ be the
$\sigma$-field of sets invariant under every permutation in $S_n$, and set
\[
\mathcal E_n=X^{-1}(\mathcal E_n'),
\qquad
\mathcal E=\bigcap_{n=1}^{\infty}\mathcal E_n.
\]
The $\sigma$-field $\mathcal E$ is the exchangeable $\sigma$-field. The tail field
\[
\mathcal T=\bigcap_{n=1}^{\infty}
 \sigma(X_{n+1},X_{n+2},\ldots)
\]
satisfies $\mathcal T\subseteq\mathcal E$.

\begin{thm}[Finite-permutation averaging]
Let $(X_n)$ be exchangeable and let
$\varphi:E^{\mathbb N}\to\mathbb R$ be integrable. For a permutation $\rho$,
write $X^\rho=(X_{\rho(1)},X_{\rho(2)},\ldots)$. For $n\geq1$, define
\[
A_n\varphi
 =\frac1{n!}\sum_{\rho\in S_n}\varphi(X^\rho).
\]
Then
\[
A_n\varphi=\mathbb E[\varphi(X)\mid\mathcal E_n].
\]
\end{thm}

\begin{proof}
The average is $\mathcal E_n$-measurable. If $A\in\mathcal E_n$, then
exchangeability and invariance of $A$ give, for every $\rho\in S_n$,
\[
\mathbb E[\varphi(X)\mathbf1_A]
 =\mathbb E[\varphi(X^\rho)\mathbf1_A].
\]
Average over $\rho$ to obtain the defining identity for conditional
expectation.
\end{proof}

\begin{thm}[Reverse martingale convergence]
Let $\mathcal G_1\supseteq\mathcal G_2\supseteq\cdots$ and
$\mathcal G_\infty=\bigcap_n\mathcal G_n$. For $Z\in L^1$,
\[
\mathbb E[Z\mid\mathcal G_n]
 \longrightarrow\mathbb E[Z\mid\mathcal G_\infty]
\quad\text{almost surely and in }L^1.
\]
More generally, if $(Y_n)$ is adapted to the decreasing filtration and
\[
\mathbb E[Y_n\mid\mathcal G_{n+1}]=Y_{n+1},
\]
then $(Y_n)$ is called a reverse martingale and has the displayed form with
$Z=Y_1$.
\end{thm}

\begin{proof}
The conditional expectations are uniformly integrable. To obtain almost sure
convergence, fix $N$ and reverse the finite sequence:
\[
Z_k^{(N)}=\mathbb E[Z\mid\mathcal G_{N-k}],
\qquad 0\leq k\leq N-1.
\]
With the increasing filtration $\mathcal G_{N-k}$, this is an ordinary
martingale. Doob's upcrossing inequality applied to this reversed martingale
and to its negative gives bounds, independent of $N$, on the expected numbers
of upcrossings and downcrossings of every rational interval by the original
sequence. Letting $N\to\infty$ and using the deterministic crossing criterion
shows that the original sequence converges almost surely. Uniform
integrability then gives $L^1$ convergence.

If $(Y_n)$ is a reverse martingale, iteration of the tower property gives
$Y_n=\mathbb E[Y_1\mid\mathcal G_n]$. Finally, testing the $L^1$ limit against
events in $\mathcal G_\infty$ identifies it as
$\mathbb E[Z\mid\mathcal G_\infty]$.
\end{proof}

\begin{thm}[Exchangeable averages and the tail field]
Let $(X_n)$ be exchangeable. Let $k\geq1$ and let
$\varphi:E^k\to\mathbb R$ satisfy
$\mathbb E|\varphi(X_1,\ldots,X_k)|<\infty$. For $n\geq k$, write
\[
A_n\varphi
 =\frac1{(n)_k}
  \sum_{\substack{1\leq i_1,\ldots,i_k\leq n\\
                  i_1,\ldots,i_k\text{ distinct}}}
  \varphi(X_{i_1},\ldots,X_{i_k}),
\]
where $(n)_k=n(n-1)\cdots(n-k+1)$. Then
\[
A_n\varphi
 \longrightarrow
 \mathbb E[\varphi(X_1,\ldots,X_k)\mid\mathcal E]
 =\mathbb E[\varphi(X_1,\ldots,X_k)\mid\mathcal T]
\]
almost surely and in $L^1$.
\end{thm}

\begin{proof}
The ordered-tuple average is the same as the permutation average, so it is
$\mathbb E[\varphi(X_1,\ldots,X_k)\mid\mathcal E_n]$. Since
$\mathcal E_n\downarrow\mathcal E$, reverse martingale convergence gives the
first limit.

Fix $\ell$. Let $B_{n,\ell}$ be the same ordered-tuple average but using only
indices in $\{\ell+1,\ldots,n\}$. It is measurable with respect to
$\sigma(X_{\ell+1},X_{\ell+2},\ldots)$. Exchangeability and a coupling of a
uniform ordered $k$-tuple from $\{1,\ldots,n\}$ with one avoiding the first
$\ell$ indices give
\[
\|A_n\varphi-B_{n,\ell}\|_1
 \leq2\mathbb E|\varphi(X_1,\ldots,X_k)|
      P(\text{the tuple uses one of }1,\ldots,\ell)
 \longrightarrow0.
\]
Thus the $L^1$ limit $L$ is measurable, in the $L^1$ sense, with respect to
every $\sigma(X_{\ell+1},X_{\ell+2},\ldots)$; equivalently,
$\mathbb E[L\mid\sigma(X_{\ell+1},\ldots)]=L$. Applying reverse martingale
convergence to $L$ shows $L=\mathbb E[L\mid\mathcal T]$, so $L$ has a tail
measurable version. This proves the equality of the two conditional
expectations.
\end{proof}

\begin{thm}[Exchangeable and tail fields agree modulo null sets]
If $(X_n)$ is exchangeable, then for every $A\in\mathcal E$ there is
$B\in\mathcal T$ such that
\[
P(A\triangle B)=0.
\]
\end{thm}

\begin{proof}
Finite-coordinate bounded functions are dense in
$L^1(\sigma(X_1,X_2,\ldots))$. Choose such functions $\varphi_m$ with
$\|\varphi_m-\mathbf1_A\|_1\to0$. The preceding theorem gives
\[
\mathbb E[\varphi_m\mid\mathcal E]
 =\mathbb E[\varphi_m\mid\mathcal T].
\]
Conditional expectation is an $L^1$ contraction, so letting $m\to\infty$
yields
\[
\mathbf1_A
 =\mathbb E[\mathbf1_A\mid\mathcal E]
 =\mathbb E[\mathbf1_A\mid\mathcal T]
\quad\text{almost surely}.
\]
Choose a tail-measurable version $Y$ of the last conditional expectation and
let $B=\{Y>1/2\}$.
\end{proof}

\begin{cor}[Hewitt--Savage zero--one law]
If $X_1,X_2,\ldots$ are i.i.d., then every exchangeable event $A$ satisfies
\[
P(A)\in\{0,1\}.
\]
\end{cor}

\begin{proof}
By the preceding theorem, $A$ agrees almost surely with a tail event.
Kolmogorov's zero--one law applies to the tail event.
\end{proof}
